\SourceLeaf{247}%
die Größe des Linienelements auf einer gekrümmten Fläche angibt. Untersuchen wir zunächst die geometrische Bedeutung der Koeffizienten \(E,F,G\). Schon in Artikel 5 haben wir darauf hingewiesen, dass man sich auf einer gekrümmten Fläche zwei Liniensysteme denken kann. Auf jeder Linie des ersten Systems ist allein \(p\) veränderlich und \(q\) konstant. Auf jeder Linie des zweiten Systems ist allein \(q\) veränderlich und \(p\) konstant. Jeder Punkt der Fläche kann als Schnittpunkt einer Linie des ersten Systems mit einer Linie des zweiten angesehen werden. Das an diesen Punkt angrenzende Element der ersten Linie, das der Änderung \(\mathrm{d}p\) entspricht, ist dann \(=\BareRoot{E}\mathord{.}\mathrm{d}p\). Das Element der zweiten Linie, das der Änderung \(\mathrm{d}q\) entspricht, ist \(=\BareRoot{G}\mathord{.}\mathrm{d}q\). Bezeichnen wir schließlich mit \(\omega\) den Winkel zwischen diesen Elementen, so sieht man leicht, dass \(\cos\omega=\frac{F}{\BareRoot{EG}}\) gilt. Der Flächeninhalt des parallelogrammförmigen Elements auf der gekrümmten Fläche zwischen den beiden Linien des ersten Systems zu \(q,q+\mathrm{d}q\) und den beiden Linien des zweiten Systems zu \(p,p+\mathrm{d}p\) ist \(\BareRoot{(EG-FF)}\mathrm{d}p\mathord{.}\mathrm{d}q\).

Eine beliebige Linie auf der gekrümmten Fläche, die zu keinem der beiden Systeme gehört, entsteht, wenn man \(p\) und \(q\) als Funktionen einer einzigen neuen Veränderlichen auffasst oder eine der beiden Größen als Funktion der anderen. Es sei \(s\) die Länge einer solchen Kurve, von einem beliebigen Anfangspunkt aus gemessen, wobei wahlweise die eine oder die andere Richtung als positiv gilt. Bezeichnen wir mit \(\theta\) den Winkel, den das Element \(\mathrm{d}s=\BareRoot{(E\mathrm{d}p^2+2F\mathrm{d}p\mathord{.}\mathrm{d}q+G\mathrm{d}q^2)}\) mit der durch seinen Anfangspunkt gezogenen Linie des ersten Systems bildet. Damit keine Mehrdeutigkeit bleibt, nehmen wir an, dass der Winkel stets von dem Zweig dieser Linie aus gemessen wird, auf dem \(p\) wächst, und nach der Seite hin positiv zählt, auf der \(q\) wächst. Mit diesen Festlegungen ergeben sich leicht die folgenden Beziehungen
\[
\begin{aligned}
\cos\theta\mathord{.}\mathrm{d}s&=\BareRoot{E}\mathord{.}\mathrm{d}p+\BareRoot{G}\mathord{.}\cos\omega\mathord{.}\mathrm{d}q=\frac{E\mathrm{d}p+F\mathrm{d}q}{\BareRoot{E}}\\
\sin\theta\mathord{.}\mathrm{d}s&=\BareRoot{G}\mathord{.}\sin\omega\mathord{.}\mathrm{d}q=\frac{\BareRoot{(EG-FF)}\mathord{.}\mathrm{d}q}{\BareRoot{E}}
\end{aligned}
\]
\Article{18}
Wir untersuchen nun, unter welcher Bedingung diese Linie die kürzeste ist. Da ihre Länge \(s\) durch das Integral
\[
s=\int\BareRoot{(E\mathrm{d}p^2+2F\mathrm{d}p\mathord{.}\mathrm{d}q+G\mathrm{d}q^2)}
\]
gegeben ist, verlangt die Bedingung für ein Minimum, dass die Variation dieses Integrals, die durch eine unendlich kleine Änderung des Linienverlaufs entsteht, \(=0\) ist. Für unseren Zweck lässt sich die Rechnung in diesem Fall leichter ausführen, wenn wir \(p\) als Funktion von \(q\) auffassen.
